\documentclass[10pt,a4paper]{amsart}
\usepackage[margin=19mm]{geometry}
\usepackage{amsmath,amssymb,mathtools}
\usepackage[scr=rsfs,cal=euler]{mathalfa}
\usepackage{xfrac}
\usepackage[unicode,hidelinks,bookmarks=false]{hyperref}
\newcommand{\ie}{i.e.\ }
\newcommand{\cf}{cf.\ }
\newcommand{\resp}{respectively\ }
\newcommand{\Subsection}{\subsection}
\providecommand{\nobreakdash}{\nobreak\hbox{-}}
\setlength{\emergencystretch}{2em}
\multlinegap0pt
\usepackage{bm}
\usepackage{bbm}
\usepackage{dsfont}
\usepackage{xspace}
\usepackage{cancel}
\usepackage{mathtools}
\usepackage{amsthm}
\makeatletter
\@ifundefined{theorem}{\newtheorem{theorem}{Theorem}}{}
\@ifundefined{lemma}{\newtheorem{lemma}[theorem]{Lemma}}{}
\makeatother
\title[Asymptotic behavior of the extremal position in a multi-type]{Asymptotic behavior of the extremal position in a multi-type branching random walk with heavy-tailed displacements}
\author{Krzysztof Kowalski}
\date{Source: arXiv:2509.14808v1 (CC BY 4.0)}
\begin{document}
\maketitle

\noindent\emph{Source and licence.} Asymptotic behavior of the extremal position in a multi-type branching random walk with heavy-tailed displacements. Version arXiv:2509.14808v1 (CC BY 4.0), distributed under the Creative Commons Attribution 4.0 International licence, as recorded in the arXiv licence metadata for this exact version and shown on the abstract page https://arxiv.org/abs/2509.14808v1. Full rights notice: licenses/arxiv-2509.14808-CC-BY-4.0.txt.\par\smallskip

\noindent\textit{Editorial scope.} Counted unit: the Lemma whose source label is \texttt{lemma3} in the article (source0.tex lines 975--986, source label \texttt{lemma3}), delivered together with the article's own complete original proof of it (source0.tex lines 987--1022, one \texttt{proof} environment of 36 lines). No mathematics is written, repaired or re-derived: statement and proof are the article's own text line for line. Required context retained uncounted: as3 (lines 127--129); lemma (lines 444--450); lemma2 (lines 452--457). Every definition, display and label of those lines is retained unchanged. Omitted from this package and not redistributed: 1--126, 130--443, 451--451, 458--974, 1023--1218. Nothing inside a retained excerpt is omitted or altered: every retained body is delivered exactly as the article's lines are written, so no omission receipt of any kind accompanies this package. Citations: the retained excerpts carry no citation command at all, so no bibliography entry of the article is transcribed and no reference is redistributed. Reference adaptation: none was needed and none was made; every reference inside the delivered unit resolves inside it, so no unresolved reference or citation is produced. Printed slips kept literal: no printed word, symbol, hypothesis or number of the counted statement or of its proof was altered. Layout and class: the article's own class is \texttt{article} with packages including amssymb, amsmath, mathtools, inputenc, hyperref, comment, mathrsfs, geometry, amsthm, amsrefs, float, fontenc, graphicx; the wrapper is the lane's pinned \texttt{amsart} editorial wrapper at 10pt on A4, which restates the theorem-like environments the retained text needs and adds the article's own macro definitions that the retained text needs (none), with extra wrapper packages (bm, bbm, dsfont, xspace, cancel, mathtools). The delivered numbering is the unit's own. Assets: no retained span contains a figure environment, a graphics-inclusion command, a PDF-page inclusion, a table or a TikZ picture; no image file is copied into this package, the package declares no asset dependency and no image file is redistributed with it. Licence: version arXiv:2509.14808v1 of this article is distributed under the Creative Commons Attribution 4.0 International licence, as recorded in the arXiv licence metadata for this exact version and shown on the abstract page https://arxiv.org/abs/2509.14808v1. A full rights notice is delivered at licenses/arxiv-2509.14808-CC-BY-4.0.txt. Characters: every retained excerpt is pure ASCII, so the delivered engine, XeLaTeX with its default Latin Modern text font, reports no missing character.
\begin{equation}\label{as3}
    M^l(i,j) > 0 \text{ for all } i,j \in \mathcal{C}.
\end{equation} 

         \begin{lemma}\label{lemma}
         Let $|Z_n| = \sum_{r \in \mathcal{C}} Z^r_n$ be the sum of all particles in the $n$-th generation of the process. Then for any $\varepsilon > 0$ there is $ 0< \delta \leq \varepsilon$ satisfying
         \[
         \mathbb{P} \left(\frac{|Z_n|} {\rho^n}  <  {(1-\varepsilon)^n} \right) < (1-\delta)^{ n}
         \]
         for all $n$ large enough. 
     \end{lemma}

     \begin{lemma}\label{lemma2}
     There exist $\delta >0$ and $\beta \in (0,1)$, such that for all $i \in \mathcal{C}$ and all $n$ large enough 
     \[
     \mathbb{P}(Z_n^i < \delta |Z_{n-l}|) \leq \beta^n,
     \]where $l \in \mathbb{N}$ is as in \eqref{as3}.
     \end{lemma}

\begin{lemma}\label{lemma3}
Denote the number of particles of class $a$ in the $n$-th generation as $Z_n(a)$. If $a \preceq j$, then, for any $\varepsilon > 0$, there exist $k \in \mathbb{N}$, $\delta >0$ and $\beta,\gamma \in (0,1)$, such that for all $n$ large enough 
     \begin{equation}\label{l3:1}
         \mathbb{P}\left(Z^j_n \leq \delta  Z_{n-k}(a)\right) \leq \beta^n.
     \end{equation}
       and
    \begin{equation}\label{l3:2}
        \mathbb{P}\left(\frac{Z_{n}^j} {\rho(a)^n}< (1-\varepsilon)^n\right) \leq \gamma^n
    \end{equation}
    
    
\end{lemma}

\begin{proof}[Proof of the Lemma \ref{lemma3}]
    We prove the lemma by induction. First observe that if $m=1$, the model reduces to the irreducible one considered in previous section and the statement of the lemma is an immediate consequence of Lemmas \ref{lemma} and \ref{lemma2}. Now assume that it holds for processes with $m$ classes and consider one with $m+1$ classes.
    Observe that the statement follows immediately from the induction assumption if $j \in \mathcal{C}_l$ for some $l\leq m$, as the last type does not contribute to the previous ones, hence we only consider the case when $j \in \mathcal{C}_{m+1}$. We start by showing \eqref{l3:1}. First note that it is trivially true if $a = m+1$, so assume $a \leq m$. Fix $i \in \mathcal{C}_a$ and $k \in \mathbb{N}$ such that $M^k(i,j) > 0$ (recall that $a \preceq j$ means that such $k$ exists for all $i \in \mathcal{C}_a$). Since \eqref{l3:1} holds for first $m$ classes, the problem can be reduced to showing existence of $\beta \in (0,1)$ satisfying
    \begin{equation}\label{l3:eq1}
       \mathbb{P}\left(Z^j_n < \delta  Z^i_{n-k}\right) \leq \beta^n.
    \end{equation}
   The proof is again similar to that of Lemma \ref{lemma2}. We denote by $Z_k^{i \rightarrow j}$ a random variable distributed as $Z_k^j$ under $\mathbb{P}_i$. Consequently, $q = \mathbb{P} \left( Z_k^{i \rightarrow j} = 0 \right) < 1$. Observe
    \begin{align*}
            \mathbb{P}\left(Z^j_n < \delta  Z^i_{n-k} \right) &= \mathbb{P} \left( \sum_{r \in \mathcal{C}} \sum_{l = 1}^{Z^r_{n-k}} Z_k^{i \rightarrow j} (l)< \delta  Z^i_{n-k} \right) \\&\leq  \mathbb{P} \left(  \sum_{l = 1}^{Z^i_{n-k}} Z_k^{i \rightarrow j} (l)< \delta  Z^i_{n-k} \right) = \mathbb{E} \left[\Phi\left(Z_{n-k}^i\right)\right]
    \end{align*}
where 
\[
\Phi(s) =  \mathbb{P} \left(  \sum_{l = 1}^{s} Z_k^{i \rightarrow j}(l) < \delta  s \right).
\]
and $\{Z_k^{i \rightarrow j} (l)\}_{l >0} $ are indpendent copies of $Z_k^{i \rightarrow j}$. 
If $q = 0$, the statement of the Lemma is trivially true with $\delta = 1$. Hence, we only consider the case where $q \in (0,1)$. Denote $K_n = \#\{l \leq s: Z_k^{i \rightarrow j}(l) > 0\}$. Then 
 \begin{align*}
     \Phi(s) = &\mathbb{P}\left(\sum_{l=1}^{s} Z_k^{i \rightarrow j}(l) < \delta s \right) = \sum_{t=0}^\infty \mathbb{P}\left(\sum_{l=1}^{s} Z_k^{i \rightarrow j}(l) < \delta s \ , \ K_n =t\right) \\&\leq \sum_{t=0}^\infty \mathbb{P}(t <\delta s,K_n=t) = \sum_{t=0}^{\lfloor \delta s \rfloor} \binom{s}{t} (1- q)^t q^{s-t} \leq q^s\sum_{t=0}^{\lfloor \delta s \rfloor}\frac {(\delta s)^t} {t!} \left(\frac {1-q} {\delta q}\right)^t
 \end{align*}
 Choosing $\delta > 0$ small enough so that \[
 \frac {1-q} {\delta q} > 1,\] and \[
 \beta_0 = q \left(\frac {1-q} {\delta q}\right)^{\delta} < 1.
 \]we get the bound 
 \[
 \Phi(s) \leq \beta_0^s.
 \]
Then 
     \begin{align*}
    \mathbb{E} \left[\Phi\left(Z_{n-k}^i\right)\right] \leq \mathbb{E}\left[\beta_0^{Z_{n-k}^i} \right] \leq \beta_0^n + \mathbb{P}\left(Z_{n-k}^i< n\right).
    \end{align*}
    Since $i \in \mathcal{C}_a$ for $a \leq m$, $\mathbb{P}\left(Z_{n-k}^i< n\right)$ decays exponentially fast by induction assumption, which ends the proof of \eqref{l3:1}. Having proved that, \eqref{l3:2} follows easily:
    \begin{align*}
        \mathbb{P}\left(\frac {Z_{n}^j} {\rho(a)}< (1-\varepsilon)^n\right) \leq \beta^n + \mathbb{P}\left(\frac {\delta Z_{n}^i} {\rho(a)}< (1-\varepsilon)^n\right)  \leq \beta^n + \gamma_0^{\frac n {\delta}} \leq \gamma^n
    \end{align*}
    for appropriate choice of $\gamma \in (0,1)$ and large enough $n$.
    \end{proof}    

\end{document}
